% Generated draft: requires mathematical and source review.
\documentclass[11pt]{article}
\usepackage[margin=25mm]{geometry}
\usepackage{fontspec}
\setmainfont{Latin Modern Roman}
\usepackage{amsmath,amssymb,mathtools,unicode-math}
\setmathfont{Latin Modern Math}
\usepackage{enumitem,needspace,xurl,hyperref}
\setlength{\parindent}{0pt}
\setlength{\parskip}{5pt}
\setlength{\emergencystretch}{3em}
\newcommand{\N}{\mathbb{N}}
\newcommand{\Z}{\mathbb{Z}}
\newcommand{\Q}{\mathbb{Q}}
\newcommand{\R}{\mathbb{R}}
\newcommand{\C}{\mathbb{C}}
\newcommand{\F}{\mathbb{F}}
\newcommand{\A}{\mathbb{A}}
\newcommand{\PP}{\mathbb{P}}
\newcommand{\E}{\mathbb{E}}
\newcommand{\eps}{\varepsilon}
\newcommand{\dd}{\,\mathrm{d}}
\begin{document}
\section*{UnsolvedMath ID: 30002861}
\textbf{Source problem code:} catalogue entry
\subsection{Definitions and mathematical statement}
The supplied record asks two interpolation questions. The following formulation uses the conventional finite, geometric interpretation of a planar triangulation; the precise admissible class of triangulations is not specified in the record and requires confirmation.

Let \(\mathcal T\) be a finite conforming collection of nondegenerate closed triangles in \(\R^2\), with distinct triangles intersecting only in a common edge, a common vertex, or the empty set. Write \(\Omega_{\mathcal T}=\bigcup_{T\in\mathcal T}T\) and let \(V(\mathcal T)\) be its vertex set. Let \(\mathcal P_d\) denote the real polynomials in two variables of total degree at most \(d\). Define \(S_d^1(\mathcal T)\) to consist of functions \(s:\Omega_{\mathcal T}\to\R\) such that \(s|_T=p_T|_T\) for some \(p_T\in\mathcal P_d\) on each triangle, and such that the values and both first partial derivatives of adjacent polynomial pieces agree on every intersection. This specifies a continuous first-order jet, including at boundary vertices, so that \(\nabla s(v)\) is well defined. No boundary values or boundary derivatives are prescribed separately.

Under this interpretation, the cubic interpolation question is whether the evaluation map
\[
E_{\mathcal T}:S_3^1(\mathcal T)\longrightarrow\R^{V(\mathcal T)},\qquad
E_{\mathcal T}(s)=(s(v))_{v\in V(\mathcal T)}
\]
is surjective for every admissible \(\mathcal T\). Equivalently, is it true that
\[
\forall\mathcal T\;\forall a:V(\mathcal T)\to\R\;\exists s\in S_3^1(\mathcal T)\;\forall v\in V(\mathcal T),\quad s(v)=a(v)?
\]

The supplied corrected formulation of the quartic question asks whether there exists an admissible \(\mathcal T\) for which the vertex-jet map
\[
J_{\mathcal T}:S_4^1(\mathcal T)\longrightarrow
\prod_{v\in V(\mathcal T)}(\R\times\R^2),\qquad
J_{\mathcal T}(s)=\bigl(s(v),\nabla s(v)\bigr)_{v\in V(\mathcal T)}
\]
is not surjective. Explicitly, the proposed failure is
\[
\exists\mathcal T\;\exists a:V(\mathcal T)\to\R\;\exists g:V(\mathcal T)\to\R^2\;\nexists s\in S_4^1(\mathcal T)
\quad\text{such that}\quad
\forall v\in V(\mathcal T),\quad s(v)=a(v),\quad\nabla s(v)=g(v).
\]
Thus failure means that some prescribed data cannot be realized on some triangulation. The original wording, ``always impossible in the general case,'' does not itself specify these quantifiers or the quartic degree; the supplied corrected statement supplies this interpretation. No assertion of generic failure, failure on every triangulation, or impossibility for every choice of data is made here. The real scalar field, finiteness, conformity, geometric interpretation, and boundary convention above are explicit formulation conventions whose agreement with the intended source must be reviewed.
\subsection{Short English statement}
With the stated triangulation conventions, the first question asks whether every assignment of real values to the vertices can be interpolated by a globally \(C^1\), piecewise cubic spline. The second asks whether some triangulation and some assignment of values and gradients admit no globally \(C^1\), piecewise quartic interpolant. The exact admissible triangulations and the original wording of the second question require source review.
\subsection{Sources}
\begin{itemize}
\item Local UnsolvedMath registry: supplied record with permanent ID \(30002861\), title ``Cubic Spline Interpolation on General Triangulations,'' exact source problem code \texttt{catalogue entry}, in the Oberwolfach Reports Open Problems collection. No individual proposer is supplied; the proposed year is recorded as 2015.
\item Supplied report citation: \emph{Multivariate Splines and Algebraic Geometry (2016)}, Oberwolfach Reports, pp.\ 1139--1200, DOI \texttt{10.4171/owr/2015/21}. \url{https://doi.org/10.4171/owr/2015/21}. The literature-source list labels this reference ``Oberwolfach Report 21.'' The supplied citation's 2016 designation and its 2015 DOI/proposed-year metadata are retained without resolving their relationship.
\item Supplied literature reference: \emph{Interpolation by C1 splines of degree at least four}. \url{https://doi.org/10.1016/S0377-0427(99)00350-7}. Authors and further bibliographic details are not supplied.
\item Literature assessment dated exactly \texttt{2026-08-21}: the supplied curation labels the problem open, describes two unresolved interpolation conjectures, and states that available results cover special triangulations or degree ranges only. This is supplied dated curation, not an independently verified assessment; time-sensitive status and cited claims require independent verification.
\item The registry labels the statement \texttt{corrected\_verified} and reports that an official report passage was checked directly and its mathematical question, definitions, and notation restored without strengthening the claim. This is supplied verification metadata, not a claim of source inspection for this draft. No exact passage, page-level problem location, or named attribution is supplied.
\end{itemize}
\end{document}
